% chuong_v3/fig_ch06_window_shape.tex
% Dạng sóng rời rạc của Cửa sổ Hanning N+1 (N = 638) và xử lý bão hòa Q1.15 ở đỉnh

\begin{tikzpicture}[>=Stealth, font=\small, scale=0.95, every node/.style={transform shape}]

    % Trục tọa độ chính
    \draw[->, thick, draw=black!70] (-0.5, 0) -- (11.8, 0) node[right, font=\scriptsize, text=black!80] {Chỉ số mẫu $n$ ($0 \dots 637$)};
    \draw[->, thick, draw=black!70] (0, -0.3) -- (0, 4.6) node[above, font=\scriptsize, text=black!80] {Hệ số lượng tử Q1.15 $w[n]$};

    % Đường biên và các mốc
    \draw[dashed, color=black!20] (0, 3.8) -- (11.2, 3.8);
    \node[left, font=\scriptsize, text=black!70] at (-0.1, 3.8) {$32.767$ (\texttt{0x7FFF})};
    \node[left, font=\scriptsize, text=black!70] at (-0.1, 1.9) {$16.384$ (\texttt{0x4000})};
    \node[left, font=\scriptsize, text=black!70] at (-0.1, 0) {$0$};

    % Đường cong bao Hanning (Màu xanh dương)
    \draw[domain=0.5:10.5, samples=100, very thick, blue!75!black] 
        plot (\x, {3.8 * (0.5 - 0.5 * cos(360 * (\x - 0.5) / 10))});

    % Trục đối xứng ở giữa (n = 318.5)
    \draw[dashdotted, thick, orange!80!black] (5.5, 0) -- (5.5, 4.1);
    \node[above, font=\scriptsize, fill=white, draw=orange!80!black, text=orange!90!black, rounded corners=2pt, inner sep=2pt] at (5.5, 4.1) 
        {Trục đối xứng ($n = 318{,}5$)};

    % Điểm mốc biên trái (n = 0)
    \filldraw[red!80!black] (0.5, 0.05) circle (2.5pt);
    \draw[thick, black!60] (0.5, 0) -- (0.5, -0.15);
    \node[below, font=\scriptsize, text=black!80] at (0.5, -0.15) {$n = 0$};
    \draw[<-, thick, red!80!black] (0.5, 0.15) -- (1.5, 1.2) 
        node[above right, font=\scriptsize, fill=white, inner sep=2.5pt, draw=red!70!black, rounded corners=2pt] 
        {\parbox{2.4cm}{\centering\textbf{Biên trái:}\\$w[0] = 1 \ne 0$}};

    % Điểm mốc biên phải (n = 637)
    \filldraw[red!80!black] (10.5, 0.05) circle (2.5pt);
    \draw[thick, black!60] (10.5, 0) -- (10.5, -0.15);
    \node[below, font=\scriptsize, text=black!80] at (10.5, -0.15) {$n = 637$};
    \draw[<-, thick, red!80!black] (10.5, 0.15) -- (9.5, 1.2) 
        node[above left, font=\scriptsize, fill=white, inner sep=2.5pt, draw=red!70!black, rounded corners=2pt] 
        {\parbox{2.4cm}{\centering\textbf{Biên phải:}\\$w[637] = 1 \ne 0$}};

    % Hai điểm đỉnh bão hòa (n = 318, n = 319)
    \filldraw[purple!80!black] (5.35, 3.8) circle (2.5pt);
    \filldraw[purple!80!black] (5.65, 3.8) circle (2.5pt);
    \draw[<-] (5.5, 3.85) -- (7.5, 3.2) 
        node[right, font=\scriptsize, fill=white, inner sep=2pt, draw=purple!80!black, rounded corners=2pt] 
        {\parbox{4.5cm}{\textbf{Đỉnh bão hòa Q1.15:}\\[1pt]
        $w[318] = w[319] = 32.767$\\[1pt]
        \emph{Kẹp $32.768 \to 32.767$ chống tràn}}};

    % Chỉ số mẫu trung tâm n = 318, 319 với hai vạch tick rõ ràng
    \draw[thick, black!60] (5.35, 0) -- (5.35, -0.12);
    \draw[thick, black!60] (5.65, 0) -- (5.65, -0.12);
    \node[below, font=\scriptsize, text=black!80] at (5.5, -0.12) {$n = 318, 319$};

    % Vùng tính LocalPeak LP (n = 212 .. 425)
    \draw[<->, thick, teal!80!black] (3.82, -0.75) -- (7.18, -0.75);
    \node[below, font=\scriptsize, text=teal!90!black, fill=white, inner sep=1.5pt] at (5.5, -0.8) 
        {\textbf{Vùng dò LocalPeak $LP$ ($n = 212 \dots 425$, 214 mẫu)}};
    \draw[dashed, color=teal!50] (3.82, 0) -- (3.82, -0.75);
    \draw[dashed, color=teal!50] (7.18, 0) -- (7.18, -0.75);

\end{tikzpicture}
